\chapter{Những Nỗi Sợ Của Học Sinh}
\markboth{Những Nỗi Sợ Của Học Sinh}{The Fears Students Carry}

\section{Vì Sao Em Sợ Gọi Tên Trước Lớp}
\begin{truyen}{Vì Sao Em Sợ Gọi Tên Trước Lớp}{Why Am I Afraid of Being Called On in Class}
\chuhoa{C}{ó những học sinh ngồi trong lớp mà tim tăng nhịp chỉ vì nghe thầy cô bắt đầu nhìn quanh chuẩn bị gọi tên.}
\textit{(Some students sit in class with their heart racing simply because the teacher begins looking around and may call on someone.)}

Các em không chỉ sợ không trả lời được.
\textit{(They are not afraid only of failing to answer.)}

Các em sợ khoảnh khắc cả lớp quay lại nhìn mình, sợ giọng mình run, sợ vài tiếng cười nhỏ, sợ cảm giác đứng một mình giữa sự chú ý.
\textit{(They fear the moment the whole class turns toward them, the shaking voice, the small laughs, the feeling of standing alone inside everyone's attention.)}

Một lần bị bêu, bị chê, hoặc bị cười có thể đủ để nỗi sợ này kéo dài rất lâu.
\textit{(One experience of being exposed, scolded, or laughed at can be enough for this fear to last a long time.)}
\end{truyen}

\begin{baihoc}
Nỗi sợ phát biểu không phải lúc nào cũng là lười. Nhiều khi nó là ký ức của một lần bị tổn thương trước đám đông.
\textit{(Fear of speaking is not always laziness. Sometimes it is the memory of having been hurt in public.)}
\end{baihoc}

\begin{ghinhoanh}
\textbf{Short English to remember:}

Public fear often grows from public hurt.

\end{ghinhoanh}

\begin{tuhocnhanh}
1. Vì sao bị gọi tên trước lớp có thể làm học sinh căng thẳng mạnh? \textit{Why can being called on create strong stress?}

2. Những tổn thương nào làm nỗi sợ này kéo dài? \textit{What hurts make this fear last?}

3. Lớp học cần thay đổi gì để gọi tên không còn là đe dọa? \textit{What needs to change so being called on stops feeling threatening?}
\end{tuhocnhanh}
\ngancach

\section{Vì Sao Em Sợ Kiểm Tra Dù Em Có Học}
\begin{truyen}{Vì Sao Em Sợ Kiểm Tra Dù Em Có Học}{Why Am I Afraid of Tests Even When I Have Studied}
\chuhoa{N}{hiều học sinh ôn rồi mà vẫn sợ. Điều đó làm các em còn hoang mang hơn: “Em có học mà sao vẫn run?”}
\textit{(Many students study and still feel afraid. That confuses them even more: “I did prepare, so why am I still shaking?”)}

Sợ kiểm tra không chỉ đến từ việc chưa học bài. Nó còn đến từ áp lực điểm số, ký ức những lần làm không tốt, nỗi sợ phụ lòng người lớn, và cảm giác một bài kiểm tra có thể định nghĩa mình.
\textit{(Fear of tests does not come only from lack of preparation. It also comes from score pressure, memories of poor performance, fear of disappointing adults, and the feeling that a test can define the self.)}

Khi ý nghĩa của bài kiểm tra bị thổi quá lớn, nỗi run sẽ mạnh hơn cả lượng kiến thức đang có.
\textit{(When the meaning of a test is inflated too much, anxiety can overpower what one actually knows.)}
\end{truyen}

\begin{baihoc}
Muốn giảm sợ kiểm tra, không chỉ cần học bài. Còn cần chữa lại cách mình gắn giá trị bản thân vào kết quả.
\textit{(Reducing test fear requires more than preparation. It also requires repairing how self-worth is tied to outcomes.)}
\end{baihoc}

\begin{ghinhoanh}
\textbf{Short English to remember:}

Test fear is often about meaning, not just content.

\end{ghinhoanh}

\begin{tuhocnhanh}
1. Vì sao học rồi vẫn có thể rất sợ kiểm tra? \textit{Why can students still fear tests even after studying?}

2. Bài kiểm tra bị thổi quá lớn theo những cách nào? \textit{In what ways is the meaning of tests often exaggerated?}

3. Cần gỡ điều gì khỏi kết quả kiểm tra? \textit{What should be untied from test results?}
\end{tuhocnhanh}
\ngancach

\section{Nếu Em Sai Trước Mọi Người Thì Sao}
\begin{truyen}{Nếu Em Sai Trước Mọi Người Thì Sao}{What If I Am Wrong in Front of Everyone}
\chuhoa{N}{ỗi sợ sai trước đám đông mạnh hơn nhiều so với nỗi sợ sai khi ở một mình.}
\textit{(The fear of being wrong in public is much stronger than being wrong alone.)}

Bởi vì khi đó, cái đau không chỉ là tri thức bị sai, mà là cảm giác mình bị nhìn thấy trong một trạng thái kém cỏi.
\textit{(Because then the pain is not only about incorrect knowledge, but about being seen in a weak state.)}

Học sinh sợ ánh mắt, sợ bạn ghi nhớ, sợ tên mình bị gắn với một lần ngốc nghếch.
\textit{(Students fear the looks, fear being remembered, fear having their name attached to one foolish moment.)}

Nếu lớp học không đủ an toàn để sai, học sinh sẽ chọn im lặng cho chắc.
\textit{(If a classroom is not safe enough for mistakes, students will choose silence for safety.)}
\end{truyen}

\begin{baihoc}
Một lớp học tốt không phải lớp ít sai, mà là lớp con người dám sai để học.
\textit{(A good classroom is not one with few mistakes, but one where people dare to be wrong in order to learn.)}
\end{baihoc}

\begin{ghinhoanh}
\textbf{Short English to remember:}

If mistakes are unsafe, silence becomes safer than learning.

\end{ghinhoanh}

\begin{tuhocnhanh}
1. Vì sao sai trước đám đông đau hơn sai một mình? \textit{Why is public error more painful than private error?}

2. Không gian học an toàn cho cái sai cần những gì? \textit{What makes a learning space safe for mistakes?}

3. Im lặng đôi khi là dấu hiệu của điều gì? \textit{Silence can be a sign of what?}
\end{tuhocnhanh}
\ngancach

\section{Em Sợ Điểm Thấp Hay Sợ Phản Ứng Sau Điểm Thấp}
\begin{truyen}{Em Sợ Điểm Thấp Hay Sợ Phản Ứng Sau Điểm Thấp}{Do I Fear the Low Score or the Reaction After It}
\chuhoa{C}{ó khi học sinh tưởng mình sợ điểm thấp, nhưng thật ra điều làm các em hoảng hơn là chuyện xảy ra sau đó.}
\textit{(Sometimes students think they fear low scores, but what truly frightens them is what comes afterward.)}

Một ánh mắt thất vọng, một câu mỉa, một bữa cơm nặng nề, một sự so sánh với người khác.
\textit{(A disappointed look, a sarcastic sentence, a heavy dinner table, a comparison with someone else.)}

Điểm chỉ là con số. Nhưng phản ứng của người lớn có thể biến nó thành một vết thương.
\textit{(A score is only a number. But adult reactions can turn it into a wound.)}

Muốn hiểu nỗi sợ học của một học sinh, đôi khi phải nhìn ra khỏi lớp học mà nhìn vào bầu không khí quanh kết quả.
\textit{(To understand a student's academic fear, we sometimes need to look beyond the classroom and into the atmosphere surrounding results.)}
\end{truyen}

\begin{baihoc}
Rất nhiều học sinh không sợ điểm trước, mà sợ cách con người đối xử với mình sau điểm.
\textit{(Many students do not fear the score first; they fear how people will treat them afterward.)}
\end{baihoc}

\begin{ghinhoanh}
\textbf{Short English to remember:}

The number hurts less than the meaning people attach to it.

\end{ghinhoanh}

\begin{tuhocnhanh}
1. Vì sao phản ứng sau điểm có thể đáng sợ hơn chính con điểm? \textit{Why can the reaction after a score be more frightening than the score itself?}

2. Bầu không khí quanh kết quả ảnh hưởng học sinh ra sao? \textit{How does the atmosphere around results affect students?}

3. Người lớn có thể giảm nỗi sợ này bằng cách nào? \textit{How can adults reduce this fear?}
\end{tuhocnhanh}
\ngancach

\section{Em Sợ Cố Rồi Vẫn Trượt}
\begin{truyen}{Em Sợ Cố Rồi Vẫn Trượt}{I Am Afraid of Trying Hard and Still Failing}
\chuhoa{Đ}{ây là nỗi sợ khiến nhiều học sinh thà không cố hết sức.}
\textit{(This is the fear that leads many students not to try fully at all.)}

Nếu không cố, các em còn giữ được một lối thoát: “Tại em chưa làm hết thôi.”
\textit{(If they do not try fully, they keep an escape route: “I did not really give my all.”)}

Nhưng nếu đã cố thật rồi mà vẫn trượt, các em sợ mình phải đối diện với điều đau hơn: có khi năng lực mình chưa đủ ở lúc này.
\textit{(But if they try seriously and still fail, they fear facing something more painful: perhaps at this moment their ability is not enough.)}

Nỗi sợ đó rất người. Nhưng nếu vì nó mà mình không dám làm thật, mình sẽ không bao giờ biết mình đứng đâu để còn xây tiếp.
\textit{(That fear is deeply human. But if it keeps us from fully trying, we never learn where we truly stand and how to build forward.)}
\end{truyen}

\begin{baihoc}
Không cố để khỏi đau chỉ làm nỗi đau kéo dài hơn. Chấp nhận thử thật là cách duy nhất để biết mình cần lớn thêm ở đâu.
\textit{(Refusing full effort to avoid pain only prolongs the pain. Accepting a real attempt is the only way to know where growth is needed.)}
\end{baihoc}

\begin{ghinhoanh}
\textbf{Short English to remember:}

Avoiding full effort delays honest growth.

\end{ghinhoanh}

\begin{tuhocnhanh}
1. Vì sao nhiều học sinh không dám cố hết sức? \textit{Why do many students not dare to try fully?}

2. “Lối thoát” mà học sinh giữ lại là gì? \textit{What escape route are they preserving?}

3. Thử thật cho mình biết điều gì? \textit{What does a real attempt reveal?}
\end{tuhocnhanh}
\ngancach

\section{Em Sợ Mình Không Đủ Thông Minh}
\begin{truyen}{Em Sợ Mình Không Đủ Thông Minh}{I Am Afraid I Am Not Smart Enough}
\chuhoa{C}{ó một nỗi sợ âm thầm đi theo rất nhiều học sinh: sợ mình thiếu một thứ gốc rễ mà người khác có sẵn.}
\textit{(A quiet fear follows many students: the fear that they lack some fundamental quality others naturally possess.)}

Khi gặp phần khó, các em không chỉ thấy bài khó. Các em thấy nó như bằng chứng rằng đầu óc mình có giới hạn thấp hơn người khác.
\textit{(When facing hard material, they do not just see difficulty. They see it as proof that their mind has a lower ceiling than others.)}

Thông minh là một khái niệm rộng và không phô bày hết qua vài bài học hay vài lần làm bài.
\textit{(Intelligence is a broad concept and does not reveal itself fully through a few lessons or a few tests.)}

Điều học sinh cần hơn là hiểu: năng lực có thể được bồi đắp nhiều qua cách học, nhịp học, sự luyện tập và môi trường.
\textit{(What students need more is this understanding: ability can be developed through methods, rhythm, practice, and environment.)}
\end{truyen}

\begin{baihoc}
Nỗi sợ “không đủ thông minh” thường lớn hơn bằng chứng thật. Đừng vội lấy khó khăn trước mắt làm bản án trọn đời.
\textit{(The fear of not being smart enough is often larger than the actual evidence. Do not turn present struggle into a lifelong verdict.)}
\end{baihoc}

\begin{ghinhoanh}
\textbf{Short English to remember:}

Present struggle is not lifelong proof of low intelligence.

\end{ghinhoanh}

\begin{tuhocnhanh}
1. Vì sao học sinh hay sợ “thiếu thông minh”? \textit{Why do students often fear lacking intelligence?}

2. Điều gì làm nỗi sợ này có vẻ rất thật? \textit{What makes this fear feel so convincing?}

3. Năng lực học có thể được bồi đắp qua những yếu tố nào? \textit{Through what factors can learning ability grow?}
\end{tuhocnhanh}
\ngancach

\section{Em Sợ Người Khác Thất Vọng Về Em}
\begin{truyen}{Em Sợ Người Khác Thất Vọng Về Em}{I Am Afraid Others Will Be Disappointed in Me}
\chuhoa{N}{hiều học sinh chịu áp lực học không phải vì yêu việc học, mà vì sợ làm ai đó thất vọng.}
\textit{(Many students carry academic pressure not from love of learning, but from fear of disappointing someone.)}

Cha mẹ kỳ vọng, thầy cô tin tưởng, người thân hay kể về em như niềm hy vọng của nhà.
\textit{(Parents expect, teachers trust, relatives speak of the student as the family's hope.)}

Ban đầu, điều đó có thể tạo động lực. Nhưng nếu kéo dài, nó biến thành nỗi nặng: em không còn học vì hiểu việc học đáng làm, mà học vì sợ ánh mắt người khác đổi khác.
\textit{(At first, that may motivate. But over time it can become a heavy burden: the student no longer studies because learning matters, but because they fear a change in others' eyes.)}

Sống bằng nỗi sợ làm người khác thất vọng khiến con người rất dễ kiệt sức.
\textit{(Living through fear of disappointing others makes a person easy to exhaust.)}
\end{truyen}

\begin{baihoc}
Kỳ vọng có thể nâng người lên, nhưng nếu đặt sai cách, nó cũng có thể đè học sinh xuống rất sâu.
\textit{(Expectation can lift a person up, but when placed wrongly, it can also press a student down deeply.)}
\end{baihoc}

\begin{ghinhoanh}
\textbf{Short English to remember:}

Expectation can inspire, but it can also crush.

\end{ghinhoanh}

\begin{tuhocnhanh}
1. Vì sao nỗi sợ làm người khác thất vọng rất mệt? \textit{Why is the fear of disappointing others so exhausting?}

2. Kỳ vọng khác áp lực ở đâu? \textit{How is expectation different from pressure?}

3. Học sinh cần nghe câu gì để được thở hơn? \textit{What sentence might help a student breathe again?}
\end{tuhocnhanh}
\ngancach

\section{Em Sợ Một Lần Trượt Sẽ Hỏng Hết Tương Lai}
\begin{truyen}{Em Sợ Một Lần Trượt Sẽ Hỏng Hết Tương Lai}{I Am Afraid One Failure Will Ruin My Whole Future}
\chuhoa{H}{ọc sinh thường nhìn một kỳ thi bằng con mắt rất tuyệt đối.}
\textit{(Students often look at one exam in an absolute way.)}

Một lần trượt, một môn thấp, một kỳ thi không như ý, và các em tưởng mọi cửa đã đóng.
\textit{(One failed exam, one low subject, one disappointing result, and they imagine every door is closed.)}

Thực tế có những cột mốc rất quan trọng thật. Không cần phủ nhận điều đó.
\textit{(It is true that some milestones matter a lot. There is no need to deny that.)}

Nhưng rất ít khi một đời người bị quyết định trọn vẹn bởi một lần vấp. Điều định hình tương lai lâu hơn là cách em phản ứng sau lần vấp đó.
\textit{(But very rarely is an entire life determined by one stumble. What shapes the future more deeply is how you respond afterward.)}
\end{truyen}

\begin{baihoc}
Một lần trượt có thể đau và đổi hướng, nhưng nó hiếm khi đủ quyền để kết thúc cả tương lai của em.
\textit{(One failure can hurt and redirect a path, but it rarely has the power to end your whole future.)}
\end{baihoc}

\begin{ghinhoanh}
\textbf{Short English to remember:}

One failure can redirect life, but it rarely defines all of it.

\end{ghinhoanh}

\begin{tuhocnhanh}
1. Vì sao học sinh hay nhìn một kỳ thi quá tuyệt đối? \textit{Why do students often see one exam too absolutely?}

2. Điều gì mới ảnh hưởng tương lai lâu dài hơn một lần vấp? \textit{What shapes the future more than one failure?}

3. Nên nói thật với học sinh về tầm quan trọng của kỳ thi như thế nào? \textit{How should we speak honestly about important exams?}
\end{tuhocnhanh}
\ngancach

\section{Đi Học Có Khi Nào Là Một Nỗi Sợ Kéo Dài Không}
\begin{truyen}{Đi Học Có Khi Nào Là Một Nỗi Sợ Kéo Dài Không}{Can Going to School Itself Become a Long Fear}
\chuhoa{C}{ó những học sinh không ghét tri thức. Các em chỉ bắt đầu sợ chính không gian đi học.}
\textit{(Some students do not hate knowledge. They begin to fear the very space of school.)}

Sợ vì bị chê, bị gọi bất ngờ, bị so sánh, bị cô lập, bị điểm thấp nhắc mãi, hoặc vì ngày nào cũng phải gồng.
\textit{(They fear being scolded, called on suddenly, compared, isolated, repeatedly reminded of poor scores, or forced to brace every day.)}

Nếu việc đi học trở thành một nỗi sợ kéo dài, vấn đề không còn nằm ở ý chí của học sinh nữa.
\textit{(If going to school becomes an ongoing fear, the issue is no longer simply the student's willpower.)}

Lúc đó người lớn phải đủ tỉnh để hỏi: nhà trường đang dạy, hay đang làm một số học sinh co rúm lại?
\textit{(Then adults must ask clearly: is the school teaching, or is it making some students shrink inward?) }
\end{truyen}

\begin{baihoc}
Khi trường học trở thành nơi gây sợ thường trực, ta cần nhìn lại môi trường, không chỉ trách học sinh yếu đuối.
\textit{(When school becomes a place of ongoing fear, we must examine the environment, not only blame students for weakness.)}
\end{baihoc}

\begin{ghinhoanh}
\textbf{Short English to remember:}

If school becomes a place of fear, the environment must be questioned.

\end{ghinhoanh}

\begin{tuhocnhanh}
1. Những yếu tố nào làm học sinh sợ chính việc đến trường? \textit{What factors make students fear going to school itself?}

2. Khi nào vấn đề không còn là chuyện “cố gắng lên”? \textit{When is the problem no longer just “try harder”?}

3. Nhà trường cần soi lại điều gì? \textit{What does the school need to examine?}
\end{tuhocnhanh}
\ngancach

\section{Nếu Em Nói Ra Mình Đang Mệt Học Có Ai Tin Không}
\begin{truyen}{Nếu Em Nói Ra Mình Đang Mệt Học Có Ai Tin Không}{If I Say I Am Exhausted by Studying, Will Anyone Believe Me}
\chuhoa{C}{ó những học sinh rất mệt nhưng không dám nói, vì sợ người lớn nghĩ mình làm quá hoặc viện cớ.}
\textit{(Some students are deeply tired but do not dare to say it, because they fear adults will think they are exaggerating or making excuses.)}

Thành ra các em im, rồi tiếp tục gồng cho tới khi học không vào, ngủ không yên, hoặc bật khóc ở những lúc rất nhỏ.
\textit{(So they stay silent and keep forcing themselves until study stops entering the mind, sleep becomes broken, or tears appear in very small moments.)}

Một nỗi mệt không được tin sẽ thành nỗi mệt cô độc.
\textit{(Exhaustion that is not believed becomes lonely exhaustion.)}

Khi học sinh hỏi bằng thái độ hoặc bằng câu nói vòng vo rằng mình có đang quá tải không, điều đầu tiên các em cần chưa chắc là lời khuyên. Có khi là một người tin rằng nỗi mệt đó là thật.
\textit{(When students ask through their attitude or indirect words whether they are overwhelmed, the first thing they need may not be advice. It may be someone who believes the exhaustion is real.)}
\end{truyen}

\begin{baihoc}
Được tin là đang mệt là bước đầu để một học sinh dám nhận trợ giúp.
\textit{(Being believed when exhausted is the first step that helps a student accept support.)}
\end{baihoc}

\begin{ghinhoanh}
\textbf{Short English to remember:}

Being believed is part of being helped.

\end{ghinhoanh}

\begin{tuhocnhanh}
1. Vì sao học sinh mệt mà vẫn không dám nói? \textit{Why do exhausted students still hesitate to speak up?}

2. Một nỗi mệt không được tin sẽ thành điều gì? \textit{What does unbelieved exhaustion turn into?}

3. Người lớn nên phản ứng thế nào khi nghe học sinh nói mình quá tải? \textit{How should adults respond when students say they are overwhelmed?}
\end{tuhocnhanh}
\ngancach
